\subsection*{Time as the Number of Motion}

Time, in Aristotle's celebrated definition, is ``the number of motion in respect of the before and after'' (\gk{ἀριθμὸς κινήσεως κατὰ τὸ πρότερον καὶ ὕστερον}). This definition, articulated at \textit{Physics} IV.11, 219b1--2, establishes time not as an independent substance---not a container within which events transpire, not an autonomous metaphysical entity---but as an aspect of motion: motion insofar as it admits of enumeration (Aristotle, \textit{Physics}, p.~60). Where there is no motion there is no time, and where motion can be numbered there is time as that numbering. Time is continuous because motion is continuous, and motion is continuous because the magnitude over which it occurs is continuous. The before and after in time correspond to the before and after in motion, which in turn correspond to the before and after in spatial magnitude. Accordingly, time inherits its continuity and its directional structure from motion, and motion inherits these from magnitude: ``the movement goes with the magnitude, and time, as we maintain, goes with motion'' (\textit{Physics} IV.11, 219a10--13). The continuity of time is therefore not a feature imposed on an otherwise discrete medium but a property derived from the continuity of spatial magnitude transmitted through the continuity of motion. Aristotle can accordingly declare that time is ``either the same thing as movement or an attribute of movement'' (Aristotle, \textit{Physics}, p.~59).

Aristotle adds a distinction that is easy to miss but load-bearing. ``Time, then, is a kind of number. Number, we must note, is used in two ways---both of what is counted or countable and also of that with which we count. Time, then, is what is counted, not that with which we count: these are different kinds of things'' (\textit{Physics} IV.11, 219b5--10). Time belongs to the first category: it is the counted number of motion, not the counting activity. This distinction will prove decisive when we turn to the soul's role, for it allows Aristotle to locate time in the structure of motion itself---as what \textit{can} be counted---while still requiring a counter for time's full actualization as numbered.

The definition is grounded experientially as well as structurally. ``We perceive movement and time together,'' Aristotle observes at \textit{Physics} IV.11, 219a4--8; ``even when it is dark and we are not being affected through the body, if any movement takes place in the mind we at once suppose that some time has indeed elapsed; and not only that but also, when some time is thought to have passed, some movement also along with it seems to have taken place.'' Neither motion nor time is epistemically prior to the other; we apprehend them together. The reciprocity is not merely epistemic but ontological---the very ontology established by the inheritance chain magnitude $\to$ motion $\to$ time---and Aristotle's phenomenological observation here is doing the double work of describing how we experience temporality and confirming the metaphysical dependence-structure that produces it.

This definition places time squarely within $A_0$. If time is the number of motion, then time cannot exist without motion; and since motion is eternal, time is eternal. Furthermore, the now (\gk{τὸ νῦν})---the instantaneous present that defines the boundary between before and after---is not a part of time but its limit, just as the point is not a part of the line but its boundary. ``The `now' determines time, insofar as time involves the before and after'' (\textit{Physics} IV.11, 219b12--13). Aristotle develops the line--point analogy carefully at \textit{Physics} IV.12, 220a5: ``Time, then, also is both made continuous by the `now' and divided at it.'' The now is at once the principle of continuity---the moving boundary that links successive durations---and the principle of division---the point at which past ends and future begins. ``What is \textit{before} is in time; for we say `before' and `after' with reference to the distance from the `now', and the `now' is the boundary of the past and the future'' (\textit{Physics} IV.14, 223a4--6). Time, then, is the continuous interval between nows, and it is through the nows that we articulate the before and after of motion.

Aristotle illuminates the now's peculiar mode of being through an experiential correlate. ``The now corresponds to the body that is carried along, as time corresponds to the motion. For it is by means of the body that is carried along that we become aware of the before and after in the motion, and if we regard these as countable we get the `now''' (\textit{Physics} IV.11, 219b22--25). The body-carried-along analogy captures what is otherwise difficult to render: the moving body is always the same body even as it traverses different positions; the now is always the same now even as it marks different moments. One in substrate, different in being. The body moves; the now marks; time is what we get when we attend to the motion as countable.

The perceptual significance of this definition emerges when we consider that numbering is itself an activity of the soul. Aristotle raises the question explicitly at \textit{Physics} IV.14, 223a21--26: could time exist without soul? If time is the number of motion, and numbering requires a counter---a soul capable of distinguishing before from after---then perhaps time depends constitutively upon soul (Aristotle, \textit{Physics}, p.~88). This is a genuine conceptual tension within the Aristotelian corpus: motion can exist without soul, for natural bodies move whether or not anyone observes them; but time qua numbered---time as the structured continuum of before and after---seems to require a soul that counts. The two-senses-of-number distinction at 219b5--10 is what makes the tension tractable. Time as \textit{what is counted} belongs to the structure of motion itself and is therefore mind-independent in its potential availability; time as \textit{fully actualized number}, however, requires the counting activity of a soul. Gibson's ecological approach to visual perception offers a modern parallel to this difficulty, insofar as Gibson insists that perception is not the passive reception of stimuli but the active pickup of information structured by the ambient array (Gibson, \textit{The Ecological Approach to Visual Perception}, p.~292). The organism does not merely undergo time; it actively structures temporal flow through its perceptual engagement with the environment.

The tension between time as mind-independent (because motion is eternal and mind-independent) and time as mind-dependent (because numbering requires a counter) is not to be resolved by collapsing one side into the other. Rather, it is precisely this tension that motivates the chain structure. $A_0$, as the joint actuality of motion and time, includes both the objective process of change and the subjective capacity to articulate that process. The soul's role as counter does not create time \textit{ex nihilo}; it actualizes the numberable aspect of motion that is already there in potentiality. Thus $A_0$ is the ground upon which the soul's constitutive contribution to time-perception rests---a contribution that becomes fully explicit at $A_1$, where the perceiving soul actively numbers the motions of the sensible world.
